\chapter{Architettura corrente}
\label{cap:sdd-architettura-corrente}

Tech4All non sostituisce alcun sistema software esistente. La prassi che
intende migliorare, descritta nel RAD, non è mediata da alcun sistema
informatico: è un insieme di procedure umane — la delega a un conoscente, la
ricerca autonoma in rete, il ricorso a un servizio di assistenza — che non ha
un'architettura da documentare.

L'assenza di un antecedente ha conseguenze precise sulla progettazione, ed è
per questo che il capitolo esiste anziché essere omesso.

\section{Vincoli che il sistema non eredita}

\begin{longtable}{@{}L{4.6cm} L{9.8cm}@{}}
\toprule
\textbf{Vincolo assente} & \textbf{Conseguenza sulla progettazione} \\
\midrule
\endfirsthead
\toprule
\textbf{Vincolo assente} & \textbf{Conseguenza sulla progettazione} \\
\midrule
\endhead
\bottomrule
\endlastfoot

Interfacce da preservare &
Nessun formato di scambio, protocollo o schema è imposto dall'esterno. La
decomposizione in sottosistemi e la forma dell'interfaccia HTTP possono essere
scelte in funzione del dominio e non della compatibilità. \\

Dati da migrare &
Non esiste una base di dati preesistente. Lo schema può essere normalizzato
liberamente e gli invarianti possono essere affidati al DBMS senza dover
accogliere dati che li violano. \\

Tecnologie imposte &
Nessuna piattaforma è vincolata da scelte pregresse. I vincoli sulla
tecnologia sono solo quelli di progetto (\id{V\_1}--\id{V\_5} del RAD). \\

Utenti da riconvertire &
Non esistono abitudini d'uso consolidate su un sistema precedente. L'unico
riferimento è la competenza digitale del pubblico, che è bassa per
definizione: è questo, e non la continuità con un sistema esistente, a
guidare le scelte di interazione. \\

\end{longtable}

\section{Vincoli che il sistema eredita comunque}

L'assenza di un sistema precedente non significa assenza di vincoli. Tre
condizioni dell'ambiente sono date e la progettazione deve accoglierle.

\begin{description}[leftmargin=1.4cm,style=nextline]
  \item[Il browser] è l'unico canale d'accesso previsto (\id{V\_1}). Il
    sistema non può presupporre installazioni sul dispositivo dell'utente né
    capacità di calcolo particolari, e deve funzionare su dispositivi datati.
  \item[La connessione] è intermittente e di qualità variabile. Il sistema è
    progettato perché ogni operazione sia autonoma: nessuna funzionalità
    richiede una sessione di lavoro continua o il completamento di più
    passaggi in sequenza entro un limite di tempo.
  \item[La competenza dell'utente] è il vincolo più stringente. Determina la
    scelta di rendere pubblica la consultazione (\id{OB\_1}), l'obbligo di
    messaggi d'errore comprensibili (\id{RNF\_U\_2}) e la richiesta di
    conferma sulle operazioni irreversibili (\id{RNF\_U\_3}).
\end{description}

Il capitolo seguente descrive l'architettura proposta a partire da queste
condizioni.
